% Bagian: turing-machines
% Bab: undecidability
% Subbagian: introduction

\documentclass[../../../include/open-logic-section]{subfiles}

\begin{document}

\olfileid[id]{tur}{und}{int}
\olsection{Pendahuluan}

Barangkali tampak jelas bahwa tidak setiap fungsi, bahkan tidak setiap fungsi
aritmetis, dapat dihitung. Jumlah fungsi terlalu banyak dan perilakunya terlalu
rumit. Misalnya, fungsi yang didefinisikan berdasarkan peluruhan partikel
radioaktif, atau berdasarkan perilaku kacau atau acak lainnya. Andaikan kita
mulai menghitung selang satu detik sejak suatu waktu tertentu, lalu
mendefinisikan fungsi $f(n)$ sebagai banyaknya partikel di alam semesta yang
meluruh dalam selang satu detik ke-$n$ setelah saat awal itu. Fungsi ini tampak
sebagai calon fungsi yang tidak mungkin dapat kita hitung.

Namun, tidak dapat membayangkan cara menghitung fungsi semacam itu merupakan
satu hal, sedangkan benar-benar membuktikan bahwa fungsi tersebut tidak dapat
dihitung merupakan hal lain. Bahkan fungsi yang tampaknya teramat rumit pun
mungkin, dalam arti abstrak, dapat dihitung. Misalnya, andaikan alam semesta
berhingga dalam waktu---pada suatu hari, jauh di masa depan, alam semesta akan
mengerut menjadi satu titik, sebagaimana diprediksi oleh beberapa teori
kosmologi. Dalam hal itu, hanya ada sejumlah berhingga (walaupun luar biasa
besar) detik sejak saat awal tersebut yang membuat $f(n)$ terdefinisi. Setiap
fungsi yang hanya terdefinisi untuk sejumlah berhingga masukan dapat dihitung:
kita dapat mencantumkan semua keluarannya dalam satu tabel besar atau
mengodekannya dalam satu diagram transisi keadaan mesin Turing yang sangat
besar.

Kita sering tertarik pada kasus khusus fungsi yang nilainya memberikan jawaban
atas pertanyaan ya/tidak. Misalnya, pertanyaan ``apakah $n$ bilangan prima?''
dikaitkan dengan fungsi
\[
\fn{isprime}(n) = \begin{cases}
  1 & \text{jika $n$ prima}\\
  0 & \text{selain itu.}
  \end{cases}
\]
Kita mengatakan bahwa suatu pertanyaan ya/tidak dapat \emph{diputuskan secara
efektif} jika fungsi terkait yang bernilai $1/0$ dapat dihitung secara efektif.

Untuk membuktikan secara matematis bahwa ada fungsi yang tidak dapat dihitung
secara efektif atau masalah yang tidak dapat diputuskan secara efektif, kita
harus menetapkan suatu model komputasi tertentu, lalu menunjukkan bahwa ada
fungsi yang tidak dapat dihitung atau masalah yang tidak dapat diputuskan oleh
model tersebut. Misalnya, kita dapat menunjukkan bahwa tidak setiap fungsi
dapat dihitung oleh mesin Turing dan tidak setiap masalah dapat diputuskan oleh
mesin Turing. Kemudian kita dapat menggunakan tesis Church--Turing untuk
menyimpulkan bahwa bukan hanya mesin Turing yang tidak cukup kuat untuk
menghitung setiap fungsi: tidak ada prosedur efektif yang dapat melakukannya.

Kunci untuk membuktikan hasil-hasil negatif semacam itu adalah kenyataan bahwa
kita dapat memasangkan bilangan dengan mesin Turing itu sendiri. Cara termudah
untuk melakukannya ialah dengan mengenumerasinya, mungkin dengan menetapkan
suatu cara khusus untuk menuliskan mesin Turing beserta programnya, lalu
mencantumkannya secara sistematis. Setelah kita melihat bahwa hal ini dapat
dilakukan, keberadaan fungsi yang tidak dapat dihitung oleh mesin Turing
mengikuti dari pertimbangan kardinalitas sederhana: himpunan fungsi dari
$\Nat$ ke~$\Nat$ (bahkan hanya dari $\Nat$ ke $\{0, 1\}$) bersifat
!!{nonenumerable}, tetapi karena kita dapat mengenumerasi semua mesin Turing,
himpunan fungsi yang dapat dihitung oleh mesin Turing hanya bersifat
!!{denumerable}.

Kita juga dapat mendefinisikan fungsi dan masalah \emph{tertentu} yang dapat
kita buktikan masing-masing tidak dapat dihitung dan tidak dapat diputuskan.
Salah satu masalah tersebut ialah apa yang disebut \emph{Masalah Penghentian}.
Mesin Turing dapat dideskripsikan secara berhingga dengan mencantumkan
instruksi-instruksinya. Deskripsi suatu mesin Turing, yakni program mesin
Turing, tentu saja dapat digunakan sebagai masukan bagi mesin Turing lain.
Karena itu, kita dapat mempertimbangkan mesin Turing yang memutuskan pertanyaan
tentang mesin Turing lain. Salah satu pertanyaan yang sangat menarik ialah:
``Apakah mesin Turing yang diberikan akhirnya berhenti ketika dijalankan pada
masukan~$n$?'' Akan berguna seandainya ada mesin Turing yang dapat memutuskan
pertanyaan ini: bayangkan mesin itu sebagai mesin Turing pengendali mutu yang
memastikan bahwa mesin Turing tidak terjebak dalam kalang tak berhingga dan
sebagainya. Fakta menarik yang dibuktikan Turing ialah bahwa mesin Turing
semacam itu tidak mungkin ada. Tidak mungkin ada satu mesin Turing yang, ketika
dijalankan pada masukan berupa deskripsi mesin Turing~$M$ dan suatu
bilangan~$n$, selalu berhenti dengan keluaran $1$ atau $0$ menurut apakah~$M$
akan berhenti ketika dijalankan pada masukan~$n$ atau tidak.

Setelah mempunyai contoh masalah tertentu yang tidak dapat diputuskan, kita
dapat menggunakannya untuk menunjukkan bahwa masalah lain juga tidak dapat
diputuskan. Misalnya, salah satu masalah terkenal yang tidak dapat diputuskan
ialah pertanyaan, ``Apakah !!{formula} orde pertama~$!A$ valid?''. Tidak ada
mesin Turing yang, apabila diberi !!{formula} orde pertama~$!A$ sebagai
masukan, dijamin berhenti dengan keluaran $1$ atau $0$ menurut apakah $!A$
valid atau tidak. Secara historis, pertanyaan untuk menemukan prosedur yang
secara efektif menyelesaikan masalah ini disebut semata-mata ``masalah
keputusan''; karena itu, kita mengatakan bahwa masalah keputusan tidak dapat
diselesaikan. Turing dan Church membuktikan hasil ini secara independen pada
waktu yang hampir bersamaan, sehingga hasil tersebut juga disebut Teorema
Church--Turing.

\end{document}
